\chapter{Коротковременное корреляционное ядро как селектор скоростей}
\label{ch:v8-correlation-kernel-short-time-rate-selector-gate}

\section{Обратная задача после конуса метрик}

Предыдущий гейт оставил шесть положительных интенсивностей
\begin{equation}
 \mathcal L_{\boldsymbol\kappa}=\sum_{r=1}^{6}\kappa_r\mathcal L_r,
 \qquad \kappa_r>0,
 \label{eq:v8-rate-selector-generator}
\end{equation}
где отношение двух ориентаций каждой пары переноса уже равно $e^{-2}$.
Проверим, содержит ли полный операторный процесс
\begin{equation}
 \mathcal C_\tau=\exp(\tau\mathcal L_{\boldsymbol\kappa})
 \label{eq:v8-rate-selector-kernel}
\end{equation}
недостающую информацию о коэффициентах.

Здесь существенно слово ``полный'': известна матрица действия
$\mathcal C_\tau$ на всей алгебре $M_{11}(\mathbb C)\oplus M_{10}(\mathbb
C)$, а не только её спектр, несколько следов или одна корреляционная
функция.

\section{Точное восстановление при калиброванном времени}

В KMS-метрике выбранный процесс подобен самосопряжённому отрицательному
оператору. Поэтому его ядро положительно и на непрерывной ветви от
$\mathcal C_0=I$ имеет однозначный логарифм:
\begin{equation}
 \mathcal L_{\boldsymbol\kappa}
 =\frac{1}{\tau}\log\mathcal C_\tau,
 \qquad
 \mathcal L_{\boldsymbol\kappa}
 =\left.\frac{d\mathcal C_\tau}{d\tau}\right|_{\tau=0}.
 \label{eq:v8-rate-selector-log-derivative}
\end{equation}
Шесть матриц $\mathcal L_r$ линейно независимы: ранг коэффициентного
дизайна равен шести, его число обусловленности равно $4.33773$. На шести
случайных положительных наборах скоростей логарифмическое восстановление
дало максимальную относительную ошибку $5.32\cdot10^{-14}$.

Если физическая длительность $\tau$ неизвестна, логарифм определяет только
$\tau\kappa_r$. Все пять независимых отношений $\kappa_r/\sum_s\kappa_s$
при этом сохраняются, но общий масштаб времени остаётся неразличимым.

Коротковременная конечная разность также сходится к генератору. При
$\tau=10^{-1}$ относительная ошибка весов равна $0.2152$, а при
$\tau=10^{-3}$ --- $0.00261$. Обратная задача, следовательно, требует либо
точного логарифма, либо действительно коротких времён.

\section{Разрешение и проверка происхождения}

Формальная обратимость не означает неограниченную численную разрешимость.
В контрольном наборе с отношением скоростей $10^4$ при $\tau=1.4$ ниже
машинной точности потеряны 163 моды: слишком быстро затухающие направления
нужно измерять раньше или с большей точностью.

Главное ограничение проекта ещё строже. Поиск по ранним формулам и Томам
VII--VIII не обнаружил независимо выведенной физической матрицы
$\mathcal C_\tau$ на концевой алгебре. Ранняя TOE-связь с тепловым ядром
имеет приближённый, постулированный характер и живёт на
пространственно-компактном носителе. Ядро предыдущего гейта Тома VIII было
синтетическим коспектральным контролем, а нынешние полугруппы строятся из
уже выбранного генератора. Использовать их затем для ``вывода'' тех же
скоростей означало бы круг.

\begin{theorem}[Условная томографическая полнота полного ядра]
На объявленном шестимерном семействе генераторов одна точная полная матрица
$\mathcal C_\tau$, принадлежащая непрерывной KMS-полугруппе и снабжённая
известным $\tau>0$, однозначно восстанавливает все шесть интенсивностей.
Без калибровки $\tau$ она однозначно восстанавливает пять относительных
интенсивностей, но не общий временной масштаб.
\end{theorem}

\begin{nogocriterion}[Запрет кругового корреляционного замыкания]
Нельзя объявлять скорости выведенными из ядра, если это ядро сначала
получено экспоненцированием генератора с теми же выбранными скоростями.
Нужен независимый родитель: действие, состояние или измеримое правило,
которое задаёт полные матричные элементы $\mathcal C_\tau$ до подстановки
$\boldsymbol\kappa$.
\end{nogocriterion}

\begin{protocol}
\begin{enumerate}
 \item независимых семейных генераторов: \textbf{6};
 \item ранг коэффициентного дизайна: \textbf{6};
 \item число обусловленности: \textbf{$4.33773$};
 \item точное восстановление при известном $\tau$: \textbf{пройдено};
 \item максимальная относительная ошибка весов: \textbf{$5.32\cdot10^{-14}$};
 \item пять относительных весов при неизвестном $\tau$: \textbf{восстановимы};
 \item общий масштаб скорости при неизвестном $\tau$: \textbf{не восстановим};
 \item коротковременная конечная разность: \textbf{сходится};
 \item барьер быстрых мод при конечной точности: \textbf{обнаружен};
 \item независимое физическое ядро концевых секторов в проекте: \textbf{не найдено};
 \item итог: \textbf{условное точное восстановление, физический источник открыт};
 \item следующий гейт: \textbf{родительское происхождение физического
 корреляционного ядра}.
\end{enumerate}
\end{protocol}

Машинный сертификат сохранён в
\path{s2t_v8_correlation_kernel_short_time_rate_selector_gate_results.json}.

\begin{thebibliography}{9}
\bibitem{v8-kernel-wolf}
M.~M.~Wolf, J.~Eisert, T.~S.~Cubitt, J.~I.~Cirac,
\emph{Assessing Non-Markovian Quantum Dynamics},
Phys. Rev. Lett. 101 (2008) 150402; arXiv:0711.3172.
\bibitem{v8-kernel-dobrynin}
D.~Dobrynin, L.~Cardarelli, M.~M\"uller, A.~Bermudez,
\emph{Compressed-sensing Lindbladian quantum tomography with trapped ions},
arXiv:2403.07462.
\bibitem{v8-kernel-liu}
Y.~Liu, J.~R.~Seddon, T.~Kohler, E.~Onorati, T.~S.~Cubitt,
\emph{Robust Lindbladian Estimation for Quantum Dynamics},
arXiv:2507.07912.
\end{thebibliography}